\subsection{局部环之间的支配关系}

定义 1. 设 A、B 是两个局部环。若 A 是 B 的子环且 \(\mathfrak{m}(\mathbf{A}) = \mathbf{A} \cap \mathfrak{m}(\mathbf{B})\)，则称 B 支配 A。

命题 1. 设 A、B 是局部环，且 A 是 B 的子环。下列条件等价：

\begin{enumerate}
\def\labelenumi{(\alph{enumi})}
\item
  \(\mathfrak{m}(\mathbf{A}) \subset \mathfrak{m}(\mathbf{B});\)
\item
  B 支配 A；
\item
  理想 \(\mathbf{B}\mathfrak{m}(\mathbf{A})\)（由 \(\mathfrak{m}(\mathbf{A})\) 在 \(\mathbf{B}\) 中生成）不包含 1。
\end{enumerate}

若 \(\mathfrak{m}(\mathbf{A}) \subset \mathfrak{m}(\mathbf{B})\)，则 \(\mathfrak{m}(\mathbf{B}) \cap \mathbf{A}\) 是 A 的一个不包含 1 且包含极大理想 \(\mathfrak{m}(\mathbf{A})\) 的理想；因此它等于该理想，即 (a) 蕴涵 (b)。若 B 支配 A，则理想 \(\mathbf{B}\mathfrak{m}(\mathbf{A})\) 含于 \(\mathfrak{m}(\mathbf{B})\)，从而不包含 1；于是 (b) 蕴涵 (c)。若 (c) 成立，则 \(\mathbf{B}\mathfrak{m}(\mathbf{A})\) 含于 B 的唯一极大理想 \(\mathfrak{m}(\mathbf{B})\)，由此得 (a)。

注意，若 K 是环，则关系“B 支配 A”是 K 的局部子环所成之集上的一个序关系。

设 A、B 是两个局部环，且 B 支配 A。典范单射 A → B 通过取商定义了从域 κ(A) 到 κ(B) 的一个子域上的同构；借助于这一同构，我们可以把 κ(A) 等同于 κ(B) 的一个子域。

\textbf{例}\par

\begin{enumerate}
\def\labelenumi{(\arabic{enumi})}
\item
  设 A 是 Noether 局部环，\(\hat{\mathbf{A}}\) 是它的完备化；则局部环 \(\hat{\mathbf{A}}\) 支配 A（第三章 § 3 第 5 节命题 9）。
\item
  设 B 是整环，A 是 B 的子环，\(p'\) 是 B 的素理想，且 \(p = A \cap p'\)。则 \(pA_{p} \subset p'B_{p'}\)，从而 \(B_{p'}\) 支配 \(A_{p}\)。
\end{enumerate}

